% \chapter{coco-block.dtx}\DescriptionDomain{Block}
%    \begin{macrocode}[numbers=none,gobble=1]
%<*block>
%    \end{macrocode}
%
% This module provides handlers for generic block elements. Blocks are
% typographical elements that are optically separated from the main
% text body, like framed boxes, quotes, transcripts, etc. This module
% provides the abstract base container for those elements.
%
% \textbf{Note:} The \lstinline{coco-block} modules diverge somewhat
% from the other {\CoCoTeX} modules insofar as that Container derived
% from Block do not follow the \CoCoTeX's usual ``collect all--process
% later'' approach, but all Properties are processed at the beginning
% of each Container's instances and the contents are processed as they
% are parsed by the \lstinline{\LaTeX} interpreter. Configuration of
% block sub-containers, however, follow the {\CoCoTeX} playbook.
%
%
% \section{Preamble}
%
%    \begin{macrocode}
%%
%% Base style for \textit{xerif} block elements.
%%
%% Maintainer: p.schulz@le-tex.de
%%
\NeedsTeXFormat{LaTeX2e}[]
\ProvidesPackage{coco-block}
    [\filedate \fileversion CoCoTeX block module]
\RequirePackage{coco-block}
%    \end{macrocode}
%
%
% \section{The Abstract Block Container}
%
% \DescribeContainer{Block} is the most abstract Container for all
% block elements.
%    \begin{macrocode}
\ccDeclareContainer{Block}{%
%    \end{macrocode}
%
%
% \subsection{Block Properties}
%
% The common properties of block elements
%    \begin{macrocode}
  \ccDeclareType{Properties}{%
%    \end{macrocode}
%
%
% \subsubsection{Margins}
%
% Margins are the spaces outside of the block element.
%
% \begin{Property}{margin-top}{<skip>} is the vertical skip at the
%   beginning of each Block instance.
%    \begin{macrocode}
    \ccSetProperty{margin-top}{\z@}%
%    \end{macrocode}
% \end{Property}
% \begin{Property}{margin-bottom}{<skip>} is the vertical skip at the
%   end of each Block instance.
%    \begin{macrocode}
    \ccSetProperty{margin-bottom}{\z@}% vertical space before the list.
%    \end{macrocode}
% \end{Property}
% \begin{Property}{margin-left}{[auto|<skip>]} is the horizontal space
%   to the left of each block instance, from the left boundary of the
%   page area. The value is passed through \lstinline{\dimexpr}, so
%   basic arithmatic is allowed.
%    \begin{macrocode}
    \ccSetProperty{margin-left}{1em}%
%    \end{macrocode}
% \end{Property}
% \begin{Property}{margin-right}{<skip>} is the right margin of the
%   list instance.
%    \begin{macrocode}
    \ccSetProperty{margin-right}{1em}% horizontal space to the right of each list item
%    \end{macrocode}
% \end{Property}
%
%
% \subsubsection{Paragraph Control}
%
% \begin{Property}{par-indent}{<skip>} is the indent of the first line
%   of a *new* paragraph inside a block instance
%    \begin{macrocode}
    \ccSetProperty{par-indent}{\parindent}%
%    \end{macrocode}
% \end{Property}
% \begin{Property}{par-fill-skip}{<skip>} is the skip at the end of
%   the last line of each paragraph inside a block instance
%    \begin{macrocode}
    \ccSetProperty{par-fill-skip}{\@flushglue}%
%    \end{macrocode}
% \end{Property}
% \begin{Property}{par-skip}{<dimen>} vertical space between two
%   adjacent paragraphs inside the block element
%    \begin{macrocode}
    \ccSetProperty{par-skip}{\z@}%
%    \end{macrocode}
% \end{Property}
%
%
% \subsubsection{Formatting}
%
% \begin{Property}{block-format}
%    \begin{macrocode}
    \ccSetProperty{block-format}{%
      \ccUseProperty{int-margin-top}%
    }%
%    \end{macrocode}
% \end{Property}
% \begin{Property}{before-block} is applied at the very beginning of
%   a~Block instance.
%    \begin{macrocode}
    \ccSetProperty{before-block}{%
      \if@noskipsec \leavevmode \fi
      \ifvmode\else
        \unskip \par
      \fi
    }
%    \end{macrocode}
% \end{Property}
% \begin{Property}{after-block} is applied at the very end of a~Block
%   instance.
%    \begin{macrocode}
    \ccSetProperty{after-block}{%
      \ccIfPropVal{after-indent}{false}{%
        \global\@afterindentfalse
        \aftergroup\cc@afterbox}{}%
    }
%    \end{macrocode}
% \end{Property}
% \begin{Property}{after-indent}{[true|false]} determines whether the
%   text paragraph after the (top-level) block element is indented
%   (\lstinline{true}) or not (\lstinline{false}).
%    \begin{macrocode}
    \ccSetProperty{after-indent}{false}%
%    \end{macrocode}
% \end{Property}
%    \begin{macrocode}
  }% /ccDeclareType{Properties}
  \ccDeclareEnv{cc@block}{endcc@block}%
}% /ccDeclareContainer{Block}
%    \end{macrocode}
%
%
%  \section{The Generic Block Environment}
%
% \begin{macro}{\cc@block} is called at the beginnning of the internal
%   block environment.
%    \begin{macrocode}
\long\def\cc@block{\cc@opt@empty\@cc@block}
\long\def\@cc@block[#1]{%
%    \end{macrocode}
% Preparing the block parameters
%    \begin{macrocode}
  \ccbCalculateVMargin{top}%
  \ccbCalculateVMargin{bottom}%
%    \end{macrocode}
% printing the block
%    \begin{macrocode}
  \ccUseProperty{before-block}%
  \ccUseProperty{int-margin-top}%
  \leftskip\dimexpr\ccUseProperty{margin-left}\relax%
  \rightskip\dimexpr\ccUseProperty{margin-right}\relax%
  \ccUseProperty{block-format}%
}
%    \end{macrocode}
% \end{macro}
% \begin{macro}{\endcc@block} is called at end of the internal block
%   environment.
%    \begin{macrocode}
\long\def\endcc@block{%
  \ccUseProperty{int-margin-bottom}%
  \ccUseProperty{after-block}
}
%    \end{macrocode}
% \end{macro}
%
%
% \section{Helper Macros}
%
% \begin{macro}{\ccbCalculateVMargin} generates a macro that sets the
%   internal margin
%   Properties\DescribeProperty{int-margin-top}\DescribeProperty{int-margin-bottom}%
%   \UsageProperty{margin-top}\UsageProperty{margin-bottom} of the
%   (nested) block element. \#1 is the orientation (\lstinline{top} or
%   \lstinline{bottom}).
%    \begin{macrocode}
\def\ccbCalculateVMargin#1{%
\message{^^J===> #1, \csmeaning{cc@\cc@cur@cont @margin-#1}}
  \ifdim\ccUseProperty{margin-#1}=\z@\relax
    \ccSetProperty{int-margin-#1}{\relax}%
  \else
    \ccSetProperty{int-margin-#1}{\addvspace{\ccUseProperty{margin-#1}}}%
  \fi
}
%    \end{macrocode}
% \end{macro}
%
%
%
%    \begin{macrocode}[numbers=none,gobble=1]
%</block>
%    \end{macrocode}
